\section{Introdução}

% Nos últimos anos, a comercialização de alimentos destinados à alimentação escolar tem ganhado relevância crescente no Brasil, especialmente no contexto do Programa Nacional de Alimentação Escolar (PNAE). Em 2025, a Lei nº 15.226 ampliou de 30\% para 45\% o percentual mínimo dos recursos do programa que devem ser destinados à aquisição direta de gêneros alimentícios da agricultura familiar, com vigência a partir de 2026~\cite{fnde_pnae_2025}. Além disso, em fevereiro de 2026, o governo federal anunciou um reajuste de 14,35\% nos repasses da merenda escolar, elevando o investimento total no programa para R\$ 6,7 bilhões~\cite{agencia_brasil_merenda_2026}. Tais medidas ampliam significativamente a importância econômica desse mercado e intensificam a participação de cooperativas, associações e pequenos produtores como agentes fornecedores.

% Nesse cenário, a questão central não se resume apenas ao processo de compra por parte do poder público, mas também às decisões estratégicas dos próprios vendedores. Cooperativas e fornecedores da agricultura familiar frequentemente ofertam múltiplos produtos (como arroz, feijão, hortaliças, frutas e derivados) para diferentes redes de ensino ou chamadas públicas, cujas condições de aquisição podem variar conforme restrições orçamentárias, custos logísticos, sazonalidade e preços de referência. Em particular, compradores distintos podem estar dispostos a aceitar valores diferentes para um mesmo item, o que torna a definição de preços uma decisão não trivial por parte do vendedor.

% Daqui pra baixo

% Relações comerciais estão presentes no cotidiano de todos. De modo geral, um vendedor expõe seus itens por um determinado valor, e um comprador, ao estar disposto a desembolsar tal quantia, pode adquiri-los. Em um mercado cada vez mais competitivo, surge uma questão central para o vendedor: como precificar seus produtos da melhor maneira possível? Nesse contexto, o problema não se limita apenas à definição de um preço, mas envolve também a compreensão do comportamento dos consumidores em diferentes cenários. Por exemplo, um comprador pode optar por adquirir um produto alternativo ao encontrar outro que ofereça a mesma funcionalidade por um preço mais acessível.

% Surge, então, um dilema econômico: preços mais elevados aumentam a receita unitária, mas podem reduzir o conjunto de compradores dispostos a adquirir determinado produto; por outro lado, preços mais baixos ampliam o número de vendas possíveis, porém diminuem o ganho obtido em cada transação. Assim, do ponto de vista do vendedor, o problema consiste em definir preços para um conjunto de itens de forma a equilibrar essas forças antagônicas e maximizar a receita total obtida.

% Esse tipo de situação se insere em uma classe mais ampla de problemas conhecidos como \textit{problemas de precificação} (PP), nos quais um vendedor deve determinar os preços de seus produtos de modo a maximizar a receita total obtida com as vendas. Essa decisão é desafiadora, pois o preço de cada item influencia simultaneamente o valor arrecadado por unidade e o conjunto de consumidores dispostos a adquiri-lo, evidenciando um \textit{trade-off} entre competitividade e rentabilidade~\cite{smith1986}.

% Problemas de precificação aparecem em diversos setores econômicos e têm sido amplamente estudados na literatura ao longo do tempo dado seu aspecto prático~\cite{smith1986,aggarwal2004,venkatesan2005,shmuel1984,shmuel1987}. Dentre as diversas variações existentes, este trabalho concentra-se no chamado \textit{Problema da Compra Mínima} (PCmin). De forma geral, nesse problema considera-se um conjunto de itens disponíveis para venda e um conjunto de consumidores, cada um com um orçamento limitado e um subconjunto de itens de interesse. O objetivo consiste em definir os preços dos itens de modo a maximizar a receita total obtida, levando em conta que cada consumidor adquire o item mais barato dentre aqueles que lhe interessam, desde que seu preço não exceda o orçamento disponível.

% Dentre os trabalhos mais recentes relacionados ao PCmin, destaca-se o de Catabriga \cite{catabriga2025}, que propõe diferentes formulações de Programação Linear Inteira (PLI) e Inteira Mista (PLIM) para o problema, além de explorar técnicas de fortalecimento dos modelos e avaliar seu desempenho computacional em diversas classes de instâncias. Seus resultados evidenciam o potencial dessas abordagens para a resolução do problema, bem como os desafios ainda existentes em termos de escalabilidade e eficiência.

% Neste contexto, o presente trabalho tem como objetivo avançar nessa linha de pesquisa, investigando possíveis melhorias nas modelagens e implementações propostas. Em particular, buscamos desenvolver formulações mais eficientes e explorar estratégias que possam reduzir o tempo de execução. Adicionalmente, pretendemos implementar e avaliar abordagens heurísticas voltadas à resolução de instâncias de grande porte, nas quais métodos exatos podem se tornar computacionalmente inviáveis, contribuindo assim para o aprimoramento das abordagens existentes na literatura.


% \blue{Daqui pra baixo eu devo falar sobre algumas trabalhos na área de pricing. Talvez iniciar falando sobre algumas variantes do problema, o que já se tem para eles (exatos, heurísticas) e levar ao artigo que define o problema formalmente. Depois disso eu devo de alguma forma falar sobre o trabalho do Sandro, falando sobre as formulações dele e o que eu pretendo fazer aqui neste trabalho}

% MARK

% Problemas de precificação consistem em determinar preços para um conjunto de produtos de modo a maximizar a receita obtida com as vendas, equilibrando um \textit{trade-off} natural entre competitividade e rentabilidade: preços mais altos aumentam o ganho por unidade, mas podem reduzir o número de compradores, enquanto preços mais baixos ampliam as vendas, porém diminuem a receita unitária~\cite{smith1986}. Entre as diversas variantes dessa classe, este projeto concentra-se no \textit{Problema da Compra Mínima} (PCmin), no qual cada consumidor, sujeito a um orçamento e a um conjunto de itens de interesse, adquire o item mais barato dentre os desejados, desde que seu preço seja viável. 

% O estudo do PCmin é relevante tanto por sua aplicação em contextos de tomada de decisão comercial quanto por sua complexidade computacional, uma vez que se trata de um problema NP-difícil. Em particular, embora trabalhos recentes, como o de Catabriga~\cite{catabriga2025}, tenham proposto formulações exatas promissoras, ainda permanecem desafios relacionados à escalabilidade dos modelos e ao tratamento eficiente de instâncias de maior porte. Nesse sentido, a principal limitação dos trabalhos correlatos está na dificuldade de manter bom desempenho computacional à medida que o tamanho das instâncias cresce.

% Diante disso, este projeto propõe investigar melhorias nas formulações matemáticas existentes para o PCmin, bem como desenvolver abordagens heurísticas para a obtenção de soluções de boa qualidade em tempo reduzido. A contribuição esperada reside tanto no fortalecimento dos modelos exatos quanto na proposição de métodos mais adequados para instâncias maiores. A validação será realizada por meio de experimentos computacionais, comparando as abordagens desenvolvidas com modelos de referência da literatura em diferentes classes de instâncias.


% Problemas de precificação consistem em determinar preços para um conjunto de produtos de modo a maximizar a receita obtida com as vendas, equilibrando um \textit{trade-off} entre competitividade e rentabilidade~\cite{smith1986}. Neste projeto, considera-se o \textit{Problema da Compra Mínima} (PCmin), no qual cada consumidor, sujeito a um orçamento e a um conjunto de itens de interesse, adquire o item mais barato dentre os desejados, desde que seu preço seja viável.

% O estudo do PCmin é relevante tanto por sua aplicação em contextos de decisão comercial quanto por sua complexidade computacional, uma vez que se trata de um problema NP-difícil. Embora trabalhos recentes, como o de Catabriga~\cite{catabriga2025}, tenham proposto formulações exatas promissoras, ainda persistem desafios relacionados à escalabilidade e à eficiência computacional em instâncias de maior porte.

% Diante disso, este projeto propõe investigar melhorias nas formulações matemáticas existentes para o PCmin e desenvolver abordagens heurísticas para a obtenção de soluções de boa qualidade em tempo reduzido. Espera-se, assim, contribuir tanto para o fortalecimento dos modelos exatos quanto para o tratamento mais eficiente de instâncias maiores. A validação será realizada por meio de experimentos computacionais, comparando as abordagens desenvolvidas com modelos de referência da literatura.


% Problemas de precificação consistem em determinar os preços de um conjunto de itens de modo a maximizar a receita total obtida com as vendas, equilibrando um \textit{trade-off} entre competitividade e rentabilidade: preços mais altos aumentam o ganho por unidade, mas tendem a reduzir o conjunto de consumidores dispostos a comprar, enquanto preços mais baixos ampliam as vendas possíveis, porém diminuem a receita unitária~\cite{smith1986}. Esses problemas possuem ampla relevância prática e têm sido amplamente estudados na literatura~\cite{smith1986,aggarwal2004,venkatesan2005,shmuel1984,shmuel1987}.

% Dentre suas diversas variantes, este trabalho concentra-se no \textit{Problema da Compra Mínima} (PCmin), no qual, dado um conjunto de itens e consumidores com orçamentos limitados e subconjuntos de interesse, busca-se definir os preços de forma a maximizar a receita total, considerando que cada consumidor adquire o item mais barato dentre aqueles que deseja, desde que seu preço seja compatível com seu orçamento.

% Entre os trabalhos recentes sobre o PCmin, destaca-se Catabriga~\cite{catabriga2025}, que propõe formulações de Programação Linear Inteira (PLI) e Inteira Mista (PLIM), além de técnicas de fortalecimento e experimentos computacionais em diferentes classes de instâncias. Embora os resultados evidenciem o potencial dessas abordagens, permanecem desafios relacionados à escalabilidade e à eficiência computacional em instâncias de maior porte.

% Nesse contexto, este trabalho busca avançar nessa linha de pesquisa por meio da investigação de melhorias nas modelagens e implementações existentes, bem como pelo desenvolvimento de abordagens heurísticas para instâncias em que métodos exatos se tornam impraticáveis, contribuindo para o aprimoramento das soluções disponíveis na literatura.


Problemas de precificação consistem em determinar os preços de um conjunto de itens de modo a maximizar a receita total obtida com as vendas, buscando um equilibrio entre competitividade e rentabilidade~\cite{smith1986}. Devido à sua relevância prática, esses problemas têm sido amplamente estudados na literatura~\cite{aggarwal2004,venkatesan2005,shmuel1984,shmuel1987,smith1986}. Dentre suas variantes, este trabalho concentra-se no \textit{Problema da Compra Mínima} (PCmin), no qual consumidores com orçamentos limitados e subconjuntos de produtos de interesse adquirem o item mais barato dentre os desejados, desde que seu preço seja viável~\cite{rusmevichientong2006}.

Entre os trabalhos recentes sobre o PCmin, destaca-se o de Catabriga~\cite{catabriga2025}, que propõe formulações de Programação Linear Inteira (PLI) e Inteira Mista (PLIM), além de técnicas de fortalecimento e análises computacionais em diferentes classes de instâncias. Embora promissoras, tais abordagens ainda apresentam desafios de escalabilidade em instâncias maiores.

Nesse contexto, este trabalho propõe investigar melhorias nas formulações existentes, bem como desenvolver abordagens heurísticas para a resolução de instâncias em que métodos exatos se tornam computacionalmente inviáveis, contribuindo para o avanço das soluções disponíveis na literatura.